\selectlanguage{russian}
Случайный поиск возникает в задачах, где агенту нужно находить редкие цели при почти полном отсутствии информации о среде: от кормодобывания животных до разведки местности, поисково-спасательных работ и поиска ресурсов. Один из ключевых открытых вопросов — какие законы движения дают наилучшие результаты при фиксированных и чётко заданных предположениях. Гипотеза о «поиске Леви» утверждает, что безмасштабное движение со степенным распределением длин перемещений может быть близким к оптимальному в двумерном пространстве. Однако выводы разных работ плохо согласуются между собой, поскольку сильно зависят от того, как именно сформулирована и реализована модель: рассматриваются ли цели как исчерпаемые или возобновляемые, прерываются ли перемещения при обнаружении, и каким образом учитываются повторные посещения.

В данной работе разработан прозрачный и воспроизводимый двумерный программный комплекс для пошагового случайного поиска на бесконечной плоскости с пуассоновским распределением целей. Модельные решения задаются явно и исключаются произвольные правила «сброса». В рамках единой постановки сравниваются четыре стратегии пост-обнаружения: продолжение полёта, телепортация и перезапуск по Висванатану в разрушающем и неразрушающем режимах. Проводится перебор показателя Леви $\mu \in [1, 3]$ (эквивалентно, хвостовой показатель $\alpha = \mu - 1 \in [0, 2]$). Эффективность оценивается главным образом по объединённой частоте уникальных обнаружений; среднее расстояние до первого обнаружения анализируется для отдельных случаев. Масштабное моделирование методом Монте-Карло для уровней разреженности $\lambda/r_v \in [10, 2000]$, дополненное анализом чувствительности параметров усечения и исследованием сходимости, позволяет построить эмпирическую карту режимов оптимального показателя $\mu^*$. Результаты согласуются с оценкой $\mu^* \approx 2$ ($\alpha \approx 1$) при наиболее разреженных протестированных условиях, показывают смещение к баллистическому движению ($\mu^* \to 1$) при высокой плотности для трёх из четырёх стратегий и документируют патологию «захвата» в неразрушающем режиме с нулевым расстоянием сброса.
